The following experiments characterize deployment beyond the native two-view, native-front-end condition without changing the CanonFuse3D architecture. Front-end adaptation uses the same fixed Unity train--validation--test split and the common 13-joint subset. VideoPose3D, PoseFormer, and MotionBERT each define two adaptation scopes: the gate and residual heads only, or all CanonFuse3D parameters. A seventh model is fully adapted on a balanced mixture of the native SAM3D stream and all three alternative front ends. Every adaptation uses the training and validation actions only, runs for 20 epochs with learning rate \(10^{-4}\), and selects the checkpoint by validation common-13 MPJPE. The frozen reference checkpoint plus seven adapted models are evaluated on all four front ends, producing an \(8\times4\) matrix of 32 test cells without test-set adaptation.

For N-view scaling, a camera group is proposed from each complete 36-camera action ring by cyclic de-duplication, yielding 180 proposed groups before data-quality filtering. Each accepted group must provide at least 30 frame indices shared by all four cameras. The nested inputs use one, two, three, and four cameras from the same group. Because CanonFuse3D is a dual-view model, the N-view prediction is formed only by composing all unique pairwise forward passes. Canonical averaging provides the non-learned baseline, the mean of pairwise CanonFuse3D predictions is the deployable aggregation, and selection of the best complete pair is reported separately as an oracle upper bound.

For the efficiency analysis, inputs are transferred to the GPU before timing. We execute ten four-view warm-up compositions, synchronize CUDA before and after each timed region, and then measure each of the 44 accepted groups once for (N=2,3,4). The timed serial path includes per-pair canonicalization, all unique pair forwards, and final arithmetic averaging; it excludes data loading, monocular front-end inference, and oracle selection. We report mean, median, and 95th-percentile latency, throughput, and peak allocated GPU memory on one NVIDIA RTX A6000 using 30-frame sequences.

The automatic synchronization experiment is an alignment-only module placed before the frozen CanonFuse3D checkpoint. For its canonical joint-velocity descriptor, every sample and view is transformed independently using that sample's own first frame; this prevents one batch element from defining another element's coordinate system. The module searches right-stream corrections from \(-5\) to \(+5\) frames in half-frame increments. Its confidence threshold is selected exclusively on the validation split by minimum offset MAE and is then frozen. Test offsets are injected at \(-5,-3,-2,-1,-0.5,0,0.5,1,2,3,\text{ and }5\) frames. All three conditions are scored on the same fixed central 20-frame interval, leaving a five-frame margin on each side, so an injected shift cannot change the evaluated pose content. The uncorrected condition uses the shifted stream, automatic uses the accepted estimated correction, and oracle restores the original unperturbed right stream. Exact lag accuracy requires absolute error below 0.25 frames; the auxiliary half-frame criterion is strict, with absolute error below 0.5 frames. A batch-order invariance test is part of the evaluation regression suite.
